\section{Complete cost accounting and the quasi-polynomial target}
\label{sec:accounting}

\subsection{A stage bound with every representation cost visible}

Fix one allocated partial complex. Let $N$ be its number of allocated
object slots, including slots made dead by previous cancellations,
$E_0$ the number of scalar differential entries, and $E_\delta$ the
number of radical entries between live objects. Put
$E_{\rm in}=E_0+E_\delta$. Let $s_1,\ldots,s_t$ be the sizes of the
weak components of the scalar graph. Write $M_{\rm sc}$ for the total
number of nonzero scalar entries retained in the sparse inclusion,
projection, and homotopy matrices. Object indices have
$O(\log(N+2))$ bits; a bound using only the live count would not
justify this assertion when many earlier slots are dead. Frontier
labels, matching identifiers, and the stored shifts contribute their
own input bit lengths.

For bit accounting, dictionary operations may be implemented by
deterministic ordered maps, with comparisons charged for their key
lengths. The resulting bounds describe the underlying algorithm;
the supplied Python implementation uses hash tables, whose timings
are reported separately. Even a conservative polynomial bound for
dictionary overhead leaves the conditional quasi-polynomial
conclusion below unchanged. We do not treat operations on arbitrarily
long Python integers as constant-time operations.

The componentwise scalar construction uses local binary matrices.
Classical elimination in a component of size $s$ has an
$O(s^3)$ binary-operation bound. The separately handled singleton and
one-edge components need constant local work each. Consequently a
conservative setup bound is
\begin{equation}
 A_{\rm sc}=
 \operatorname{poly}(\log(N+2)+\ell)\left(
 N+E_{\rm in}+M_{\rm sc}+\sum_{j:\,s_j>2}s_j^3\right),
 \label{eq:scalar-local-cost}
\end{equation}
where $\ell$ bounds the bit lengths of the non-index labels used by the
implementation. The $E_{\rm in}$ term is necessary: the construction
examines every differential entry to discover which ones are scalar.
Only a separately supplied scalar-edge index would permit charging
$E_0$ alone for that discovery. The displayed factor allows sorting
and integer-index operations; it is not intended as a
tight bit-level implementation model. The global predecessor can
instead pay cubic work in a whole matching/shift block even when that
block is disconnected, and stores global packed bit columns.

Across all adjacent homological degrees, the transfer graphs have
\[
 V=O(N+\mu),\qquad E=O(E_\delta+M_{\rm sc}),
\]
where $\mu\leq N$ is the number of scalar survivors. Each original
object appears in at most two adjacent-degree graphs, and each
relevant scalar matrix entry appears a bounded number of times. This
counts the graph representation, not only the retained corridors.

For a retained weak component $G_j$, let $V_j,E_j$ be its numbers of
vertices and edges, and let $a_j,b_j$ be its input and output counts.
The inexpensive endpoint policy chooses
\[
 r_j=\min(a_j,b_j)
\]
and evaluates from the corresponding side. It uses Boolean
reachability flags, not one packed endpoint set per vertex. Let
$\Lambda=\operatorname{poly}(\log(N+2)+\ell)$ absorb index,
dictionary, and ordering overhead. If $C(k,\ell)$ bounds the
arithmetic bit cost of a coefficient addition or typed composition
at a frontier of $2k$ endpoints, a full conservative stage bound is
\begin{equation}
 \begin{aligned}
 T_{\rm stage}\ &\leq A_{\rm sc}+A_{\rm graph}
   +\Lambda\sum_j V_jr_j\\
 &\quad+\Lambda(1+C(k,\ell))\sum_j E_jr_j.
 \end{aligned}
 \label{eq:full-stage}
\end{equation}
Here $A_{\rm graph}$ includes reading all initial graph edges,
topological ordering, pruning, and connected-component construction.
For the Boolean endpoint policy they admit a bound polynomial in
the index and label lengths times $N+V+E$. Initializing the $\mu$
output rows is included; writing a nonzero output coefficient is
charged to its terminal propagation, up to a constant factor. The
$V_jr_j$ term is required because the implementation visits the
component's ordered vertex list once for each chosen endpoint, even
when a particular coefficient is zero.

The finer reach-set policy replaces the edge sum by
$\sum_j\min(F_j,B_j)$ with unit edge weights, as proved in the
preceding section. In the vertex term, $r_j$ must then be the
number of endpoints on the \emph{actually chosen side}: a smaller
reach-set bound need not choose the side with fewer endpoints.
It must additionally pay for the packed reachability sets, with
a conservative $O((V+E)\mu)$ bit-operation term, apart from index
and label overhead. In particular, an OR or a population count on
a $\mu$-bit endpoint set is not one bit operation. This policy's
tighter propagation count can be outweighed by its larger setup cost.

With the existing dense coefficient engine, one may take
$C(k,\ell)=\operatorname{poly}(k,\ell)2^{O(k)}$. A claim about a
particular software configuration must use the coefficient engine
actually selected; the sparse default may inspect more monomial
pairs. Equation~\eqref{eq:full-stage} remains valid with the
corresponding conservative composition bound.

\subsection{Why the sparse refinement matters}

Even a column containing a single bit at global position $j$ takes
$j+1$ bits when stored as the integer $2^j$. Thus $R$ unrelated
singleton inclusions stored at positions $0,\ldots,R-1$ use
\[
 \sum_{j=0}^{R-1}(j+1)=R(R+1)/2
\]
bits of packed integer payload. Their sparse index lists use
$O(R\log(N+2))$ bits. This is a representation effect, independent
of the mathematical sparsity of the contraction.

On the shared-suffix family proved above, every scalar component is
a singleton or one identity pair. The componentwise scalar setup has
linear nonzero count, and the Boolean endpoint policy avoids the
quadratic packed reach-set volume. Reverse transfer then has
$R+2M$ radical-edge propagations, whereas the inherited forward
transfer has $R(1+2M)$. Reading the input and writing the $R$ output
coefficients are still required.

More precisely, the complete-stage separation theorem applies to
the explicitly constructed family with $M$ odd, the fixed
eight-dimensional dot algebra on one three-arc matching, sparse
scalar components, Boolean corridor flags, and the automatic
endpoint-count policy. There are $R+2M+3$ objects and $O(R+M)$
input entries, scalar supports, graph vertices, and graph edges.
For $R>1$ its one output endpoint makes the policy choose reverse
evaluation, which visits the graph once; for $R=1$ either direction
has the same bound. With generic ordering, the complete reduction uses
$O((R+M)\log(R+M+2))$ word operations and $O(R+M)$ words.
Words here contain the object indices and labels of this family;
their size is $O(\log(R+M+2))$ bits. A conservative bit implementation
therefore has a near-linear bound with additional logarithmic
factors, rather than a claimed linear bit cost.

For $R=M$ odd, inherited forward propagation alone performs
$\Omega(M^2)$ operations. The fixed coefficient algebra is essential
to this comparison: coefficient sizes are constant in this family.
This establishes the stated stage separation while including
scalar setup, and does not imply a linear running time for arbitrary
knot scans or superiority over ordinary minimum-fill cancellation.

\subsection{What still prevents a general scanner bound}

The current scanner explicitly allocates its objects before reducing
them. Neither the number of matching types nor a small boundary
alone bounds their multiplicities. The new reducer cannot avoid
reading an explicitly supplied complex of size $N$, and its minimal
output may itself contain exponentially many objects. Choosing better
scalar bases may reduce fill and corridor sizes; it does not change
the graded scalar homology dimensions of a fixed prefix.

For example, consider diagrams with normalized edge labels, charging
normalization of a longer external encoding at its actual input
bit cost. Suppose an order is supplied or constructed within the
target quasi-polynomial budget, and an execution satisfies
$N_i\leq\exp(O(\log^2(n+2)))$ and
$k_i=O(\log^2(n+2))$ at every stage, with label and shift bit lengths
polynomial in $n$. Here $N_i$ bounds each allocated stage before
reduction, including retained dead slots. Then $E_{{\rm in},i}$ and
$M_{{\rm sc},i}$ are polynomial in $N_i$, as are all graph and
endpoint counts. The setup bounds, coefficient envelope, and the
ordinary polynomial-cost crossing-extension operations on these
explicit tables make the complete $n$-stage scan quasi-polynomial.
This is sufficient accounting under explicit execution and
encoding hypotheses. There is no proof here that arbitrary diagrams
admit orders satisfying those hypotheses. The remaining research task
is to prove a structural bound or devise a representation that avoids
the explicit-object lower bound.

The same caution applies to the corridor support. Its form depends
on the chosen scalar contraction, so it is not a knot invariant and
is not canonical under arbitrary scalar basis changes. A dense basis
change can enlarge inclusion, projection, and homotopy supports while
leaving the homotopy type unchanged.

\subsection{What the covering theorem supplies to a hierarchy}

Suppose a compressed cutting step exposes a parallelity-bundle base
as a valid dihedral cover of a smaller bordered surface. The
three-family theorem then supplies the base topology, the orientation
interval-bundle type, and marked boundary lookup in time polynomial
in the presentation length and the bit length of its sheet count.
This removes a possible dependence on the expanded number of sheets
for that promised input class.

The general compressed cutting theorem of Lackenby is broader and is
already stated with logarithmic dependence on normal-surface weight
\cite[Theorem 14.4]{Lackenby2026}. Our contribution is an explicit,
small, checked kernel and classification for a restricted arithmetic
presentation. It should not be described as a new proof of the
general theorem. External attaching maps remain indispensable: two
components with the same intrinsic topology can attach differently
to the nonparallel part of a three-manifold.

\subsection{Hierarchy depth, restarts, and conditional arithmetic}

Lackenby's Section 9 yields the iteration bound
\begin{equation}
 L(g+1)^L,
 \label{eq:lackenby-iterations}
\end{equation}
where $L$ bounds hierarchy length and $g$ bounds decomposing-surface
pattern complexity \cite[Proposition 9.1]{Lackenby2026}. The integer
potential responsible for the bound is a base-$(g+1)$ encoding of the
successive complexities. A strict decrease proves termination; the
size of the integer still determines the worst possible number of
restarts. Section 10 obtains a linear bound on $L$ in the initial
handle complexity after additional parallelity-bundle operations.
Linear depth alone does not make \eqref{eq:lackenby-iterations}
quasi-polynomial.

The polynomial simplification statement of Proposition 14.2 assumes
an irreducible manifold with essential boundary pattern. Its use
inside a recognizer must respect that hypothesis. A short certificate
for an essential hierarchy also does not give a deterministic search
procedure with the same cost.

A prospective complete algorithm would be quasi-polynomial if it
visited at most $\exp(O(\log^2 n))$ compressed states, each state had
polynomial or quasi-polynomial bit size, and each required transition
had a correspondingly bounded verified implementation. The two
kernels reduce specific transition or representation costs. They do
not prove the state-count hypothesis. This is the exact remaining
gap between the delivered results and the user's global objective.
